\chapter{The Coupler and Surface Fluxes}\label{chap:ch13_coupler_fluxes}

% Local symbols (kept minimal; not provided by the preamble).
\providecommand{\Cd}{C_{\!D}}        % bulk drag coefficient
\providecommand{\Ch}{C_{\!H}}        % bulk heat-exchange coefficient
\providecommand{\Ce}{C_{\!E}}        % bulk moisture-exchange coefficient
\providecommand{\Lvap}{L_v}          % latent heat of vaporization (chapter-local; \Lvap is taken by ch09)
\providecommand{\qsat}{q_{\mathrm{sat}}}
\providecommand{\Qnet}{Q_{\mathrm{net}}}
\providecommand{\tausfc}{\boldsymbol{\tau}}

A coupled Earth System Model is, at bottom, a contract between subsystems that
otherwise know nothing of one another. The atmosphere does not know what the sea
surface temperature (SST) is; the ocean does not know how hard the wind is
blowing. The \emph{coupler} is the layer that brokers this exchange: it moves the
ocean's surface state up to the atmosphere, runs the atmosphere, computes the
fluxes of heat, freshwater and momentum across the air--sea interface from bulk
formulae, and hands those fluxes back down to drive the ocean. This chapter
documents how \model{} does this, derived directly from the source.

Two facts about the \model{} coupler are worth stating up front, because they
shape everything below. First, the \emph{active} coupled path is the multi-level
\dino{} atmosphere coupled to the $z$-level ocean
(Chapters~\ref{chap:ch04_ocean_equations} and~\ref{chap:ch09_atmosphere_multilevel}),
expressed as a single differentiable \jax{} function so that the whole coupled
step flows through \code{jax.jit}, \code{jax.grad} and \code{jax.lax.scan}
\citep{jax2018}. Second, the coupler does \emph{not} conserve energy by
construction; it uses a \emph{flux correction} to hold the mean climate against
model drift, and a substantial part of this chapter is an honest accounting of
what that correction does and what it costs.

\section{The coupling cycle}

A single coupling interval (default one model day) realizes the data flow
\begin{equation}
  \mathrm{SST}^{\,\mathrm{lin}}
  \xrightarrow{\text{regrid}} \mathrm{SST}^{\,\mathrm{Gauss}}
  \xrightarrow{\dino{}} \text{surface fields}^{\,\mathrm{Gauss}}
  \xrightarrow{\text{regrid}} \text{bulk fluxes}^{\,\mathrm{lin}}
  \xrightarrow{\text{ocean scan}} \text{ocean}(t{+}\Delta t),
  \label{eq:cycle}
\end{equation}
where ``lin'' is the regular latitude--longitude (ocean/land) grid and ``Gauss''
is the \dino{} Gaussian spectral-transform grid. The lower boundary the
atmosphere actually sees is not the bare SST but a surface skin temperature that
blends ocean, sea ice and land,
\begin{equation}
  T_{\mathrm{surf}} =
  \begin{cases}
    (1-A)\,T_{\mathrm{sst}} + A\,(T_{\mathrm{ice}} + 273.15) & \text{over ocean},\\[2pt]
    T_{\mathrm{land}} & \text{over land},
  \end{cases}
  \label{eq:lowerbc}
\end{equation}
with $A$ the (lagged) sea-ice concentration. Equation~\eqref{eq:lowerbc} uses the
\emph{previous} interval's ice and land skin temperatures: the components are
coupled \emph{explicitly} (sequentially), not implicitly, so each subsystem is
forced by its partners' state at the start of the interval.
\codref{chronos\_esm/coupler/dino\_step.py}{\code{DinoCoupledModel.\_advance} assembles $T_{\mathrm{surf}}$ and drives the cycle of Eq.~\eqref{eq:cycle}}

\section{Bulk aerodynamic surface fluxes}

The turbulent exchange of heat, moisture and momentum through the atmospheric
surface layer is not resolved at T31; it is \emph{parameterized} by bulk
aerodynamic formulae \citep{vallis2017,largeyeager2004}. These express each surface flux as the
product of a transfer coefficient, the near-surface wind speed, and an air--sea
contrast in the transported quantity.

Let $\uvec_a=(\uu_a,\vv_a)$ be the atmospheric bottom-level wind, $T_a$ and $q_a$
the bottom-level air temperature and specific humidity, and $T_s$ the surface
(SST) temperature. The effective surface wind speed used in the fluxes is floored
to represent sub-grid gustiness and capped to suppress a wind--evaporation
feedback runaway (WISHE),
\begin{equation}
  |\uvec_a|_{\mathrm{eff}} = \min\!\big(\max(\sqrt{\uu_a^2+\vv_a^2},\,1),\,25\big)\ \mathrm{m\,s^{-1}}.
  \label{eq:windeff}
\end{equation}
The sensible and latent heat fluxes (positive \emph{upward}, out of the ocean) are
\begin{align}
  Q_{\mathrm{sens}} &= \dens_a\,c_{p,a}\,\Ch\,|\uvec_a|_{\mathrm{eff}}\,(T_s - T_a),
  \label{eq:sens}\\
  Q_{\mathrm{lat}}  &= \dens_a\,\Lvap\,\Ce\,|\uvec_a|_{\mathrm{eff}}\,\big(\qsat(T_s) - q_a\big)\,\beta,
  \label{eq:lat}
\end{align}
with $\dens_a=1.2\ \mathrm{kg\,m^{-3}}$, $c_{p,a}=1004\ \mathrm{J\,kg^{-1}\,K^{-1}}$,
$\Lvap=2.5\times10^6\ \mathrm{J\,kg^{-1}}$, transfer coefficients
$\Ch=\Ce=1.2\times10^{-3}$, and an evaporation efficiency $\beta=1$ over the
ocean. The saturation specific humidity follows the Clausius--Clapeyron form
\begin{equation}
  \qsat(T_s) = \frac{0.622\,e_{\mathrm{sat}}(T_s)}{p_0},
  \qquad
  e_{\mathrm{sat}}(T_s) = 611\,\exp\!\left(\frac{17.27\,(T_s-273.15)}{(T_s-273.15)+237.3}\right)\ \mathrm{Pa},
  \label{eq:qsat}
\end{equation}
with a fixed reference $p_0=101325\ \mathrm{Pa}$. The fixed $p_0$ is an honest
simplification: it ignores the (modest) dependence of $\qsat$ on actual surface
pressure, an error of at most a few percent away from sea level.
\codref{chronos\_esm/atmos/physics.py}{\code{compute\_surface\_fluxes} implements Eqs.~\eqref{eq:windeff}--\eqref{eq:qsat}}

\paragraph{Net surface heat flux.} The bulk turbulent fluxes are combined with
radiation to form the net heat flux into the ocean,
\begin{equation}
  \Qnet = \underbrace{(1-\alpha_o)\,S^{\downarrow}}_{\text{shortwave in}}
        + \underbrace{0.8\,\sigma\,\max(T_a-10,150)^4}_{\text{longwave down}}
        - \underbrace{0.98\,\sigma\,T_s^4}_{\text{longwave up}}
        - Q_{\mathrm{sens}} - Q_{\mathrm{lat}},
  \label{eq:qnet}
\end{equation}
where $\sigma=5.67\times10^{-8}\ \mathrm{W\,m^{-2}\,K^{-4}}$ is the
Stefan--Boltzmann constant, $\alpha_o$ is the ocean albedo, and $S^{\downarrow}$
is the daily-mean top-of-atmosphere insolation reduced for cloud/atmospheric
absorption (a fixed perpetual day-of-year $80$ in the coupled harness, so there is
\emph{no seasonal cycle} in the shortwave forcing --- a known limitation). The
downwelling longwave uses a fixed effective emissivity of $0.8$ and a floored air
temperature; the upwelling term uses ocean emissivity $0.98$. This is a grey,
single-band radiation closure --- adequate to set the surface energy balance, but
not a spectral radiative-transfer scheme.
\codref{chronos\_esm/coupler/dino\_coupling.py}{\code{ocean\_fluxes\_jax} assembles $\Qnet$ from Eq.~\eqref{eq:qnet}}

\paragraph{Freshwater flux.} Evaporation is recovered from the latent heat flux,
$E = Q_{\mathrm{lat}}/\Lvap$, and the surface freshwater flux (precipitation minus
evaporation, the ocean's salinity forcing of Chapter~\ref{chap:ch04_ocean_equations}) is
\begin{equation}
  F_w = P - E = P - Q_{\mathrm{lat}}/\Lvap,
  \label{eq:fw}
\end{equation}
with $P$ the precipitation diagnosed by the atmosphere's large-scale
condensation. Both $\Qnet$ and $F_w$ are subsequently constrained for global
balance (Section~\ref{sec:fluxcorr}).

\section{Wind stress}

Momentum is transferred to the ocean by the surface wind stress
$\tausfc=(\tau_x,\tau_y)$, computed from the same bulk-drag idea but with its own
quadratic form,
\begin{equation}
  \tau_x = \dens_a\,\Cd\,|\uvec_a|_{w}\,\uu_a,
  \qquad
  \tau_y = \dens_a\,\Cd\,|\uvec_a|_{w}\,\vv_a,
  \qquad
  |\uvec_a|_{w} = \max\!\big(\sqrt{\uu_a^2+\vv_a^2},\,1\big),
  \label{eq:stress}
\end{equation}
with drag coefficient $\Cd=1.3\times10^{-3}$. Note that the wind speed floor here
($1\ \mathrm{m\,s^{-1}}$) matches the heat-flux floor but there is \emph{no}
$25\ \mathrm{m\,s^{-1}}$ cap; instead each stress component is hard-clamped,
\begin{equation}
  \tau_x,\tau_y \in [-0.3,\,0.3]\ \mathrm{Pa}.
  \label{eq:stressclamp}
\end{equation}
The clamp is a numerical safety net: realistic stresses rarely exceed
$\sim\!0.3\ \mathrm{Pa}$, so saturating it flags an instability rather than
shaping the climatology. This stress is the boundary forcing of the ocean
momentum equation (Chapter~\ref{chap:ch07_ocean_prognostic_core}) and drives the
wind-driven gyres and, indirectly, the overturning of
Chapter~\ref{chap:ch08_ocean_overturning}.
\codref{chronos\_esm/coupler/dino\_coupling.py}{\code{ocean\_fluxes\_jax}, Eqs.~\eqref{eq:stress}--\eqref{eq:stressclamp}}

\section{Flux correction: control mode vs.\ free mode}\label{sec:fluxcorr}

A coupled model with imperfect components drifts: small biases in the bulk fluxes
accumulate into large SST errors over years. \model{} controls this with a
surface-heat \emph{flux correction} on the SST, and it offers two distinct modes
selected at construction time. This is the single most important design choice in
the coupler, and it has direct consequences for forcing experiments
(Chapter~\ref{chap:ch14_forcing_response}).

\subsection{Strong Haney restoring (control mode)}

In control mode the net heat flux receives a Haney restoring term \citep{haney1971} toward a fixed
target SST field $T^{\star}$ (the World Ocean Atlas climatology captured at
initialization),
\begin{equation}
  Q_{\mathrm{corr}} = \lambda\,(T^{\star} - T_s),
  \qquad
  \lambda = \frac{\densz\,c_{p,w}\,\Delta z_1}{\tau},
  \label{eq:haney}
\end{equation}
where $\lambda$ is the restoring coefficient, $\densz=1025\ \mathrm{kg\,m^{-3}}$
and $c_{p,w}=3985\ \mathrm{J\,kg^{-1}\,K^{-1}}$ are seawater density and heat
capacity, $\Delta z_1$ is the top model-layer thickness, and $\tau$ is the
restoring timescale. With the default $\tau=30\ \mathrm{d}$ this gives
$\lambda\approx50\ \mathrm{W\,m^{-2}\,K^{-1}}$ --- a \emph{strong} relaxation that
pins the surface to climatology within roughly a month. The corrected net flux
that drives the ocean is $\Qnet + Q_{\mathrm{corr}}$.
\codref{chronos\_esm/ocean/flux\_correction.py}{\code{restoring\_lambda}, \code{restoring\_flux}, \code{heat\_correction} ($\code{q\_flux}{=}\code{None}$)}

\begin{remark}
The strong restoring is what makes the control climatology look good, but it does
so by \emph{absorbing} any imposed forcing: if you add CO$_2$ forcing while
$\lambda\approx50\ \mathrm{W\,m^{-2}\,K^{-1}}$ is acting, an SST anomaly of order
$\Delta T = F/\lambda$ develops --- a $3.7\ \mathrm{W\,m^{-2}}$ forcing yields only
$\sim\!0.07\ \mathrm{K}$. Control mode is therefore the wrong tool for climate
sensitivity; it is the right tool for spin-up and for diagnosing the implied
correction below.
\end{remark}

\subsection{Frozen q-flux + weak anomaly restoring (free mode)}

Free mode breaks the suppression without re-opening the cold-tongue drift. The
recipe is the classic mixed-layer flux correction:
\begin{enumerate}
  \item Run a strongly-restored control to equilibrium and accumulate the
        time-mean restoring flux $\bar{Q} = \langle \lambda\,(T^{\star}-T_s)\rangle$.
        This is the \emph{implied} flux correction the restoring was silently
        supplying to hold the mean state.
  \item In the free run, \emph{prescribe} the frozen field $\bar{Q}$ plus only a
        \emph{weak}, long-$\tau$ anomaly restoring:
        \begin{equation}
          Q_{\mathrm{corr}} = \underbrace{\lambda_{\mathrm{weak}}\,(T^{\star}-T_s)}_{\text{weak anomaly restoring}} + \underbrace{\bar{Q}}_{\text{frozen q-flux}},
          \qquad \tau \sim 3650\ \mathrm{d}.
          \label{eq:qflux}
        \end{equation}
\end{enumerate}
The mean state is now held by the fixed field $\bar{Q}$, residual drift is bounded
by the weak term, but SST is free to respond to an imposed forcing on the slow
anomaly timescale. Crucially the entire correction is pure \jax{}, so
$\mathrm{d}(\text{climate})/\mathrm{d}(\text{forcing})$ survives it --- the
differentiability the model is built for is not severed by the flux correction.
\codref{chronos\_esm/ocean/flux\_correction.py}{\code{heat\_correction} with \code{q\_flux} provided, Eq.~\eqref{eq:qflux}}

\begin{example}[Selecting the mode]
A \code{DinoCoupledModel} constructed with \code{q\_flux=None, restore\_tau\_days=30}
runs in control mode; passing \code{q\_flux=<field>, restore\_tau\_days=3650}
switches to free mode. The same call also balances the \emph{uncorrected} branch:
if no target is given, $\Qnet$ has its area-weighted ocean mean removed so the
global ocean neither gains nor loses heat,
\begin{equation}
  \Qnet \leftarrow \Qnet - \frac{\sum_{\text{ocean}} w\,\Qnet}{\sum_{\text{ocean}} w},
  \qquad w = \cos\lat,
\end{equation}
and the freshwater flux $F_w$ is de-meaned identically so it neither adds nor
removes net mass. Finally $\Qnet$ is clipped to $[-1500,1500]\ \mathrm{W\,m^{-2}}$
as a stability guard. \codref{chronos\_esm/coupler/dino\_coupling.py}{\code{ocean\_fluxes\_jax} mean-removal and clip}
\end{example}

\section{CO$_2$ radiative forcing into the ocean heat budget}\label{sec:co2}

The model's single anthropogenic forcing channel is CO$_2$, entered through the
Myhre logarithmic radiative-forcing law \citep{myhre1998},
\begin{equation}
  F_{\mathrm{CO_2}}(C) = 5.35\,\ln\!\left(\frac{C}{C_0}\right)\ \mathrm{W\,m^{-2}},
  \qquad C_0 = 280\ \mathrm{ppm},
  \label{eq:myhre}
\end{equation}
which vanishes at the pre-industrial reference $C_0$ and gives the canonical
$\approx3.7\ \mathrm{W\,m^{-2}}$ at doubling. A deliberate design decision is
\emph{where} this enters: $F_{\mathrm{CO_2}}$ is added to the ocean surface heat
budget as extra downwelling longwave, \emph{before} the flux correction, rather
than as an offset to the SST-slaved atmospheric radiative-equilibrium temperature.
The latter would be nearly inert (the atmosphere is slaved to SST) and would
double-count. Concretely, Eq.~\eqref{eq:qnet} becomes
$\Qnet \leftarrow \Qnet + F_{\mathrm{CO_2}}(C)$ and only \emph{then} receives
$Q_{\mathrm{corr}}$. Because $F_{\mathrm{CO_2}}$ is differentiable in $C$, the
sensitivity $\mathrm{d}(\text{climate})/\mathrm{d}(\mathrm{CO_2})$ is directly
available for calibration and sensitivity studies (Chapter~\ref{chap:ch14_forcing_response}).
\codref{chronos\_esm/coupler/dino\_coupling.py}{\code{co2\_forcing\_wm2} (Eq.~\eqref{eq:myhre}) added in \code{ocean\_fluxes\_jax} before the flux correction}

\begin{remark}
Putting $F_{\mathrm{CO_2}}$ before the correction makes the mode choice
load-bearing: under the strong control-mode $\lambda\approx50\ \mathrm{W\,m^{-2}\,K^{-1}}$
the response is suppressed to $\sim\!F/\lambda$; only free mode
(Eq.~\eqref{eq:qflux}) lets $F_{\mathrm{CO_2}}$ drive a real warming. The
historical forcing record (CO$_2$ in ppm, plus a small TSI cycle) is supplied by
\code{ForcingManager}, which linearly interpolates annual-mean concentrations
between 1850 and 2025. \codref{chronos\_esm/forcing.py}{\code{ForcingManager.get\_forcing} interpolates $C(t)$ used in Eq.~\eqref{eq:myhre}}
\end{remark}

\section{Regridding between atmosphere and ocean}

The atmosphere and ocean live on different horizontal grids: the \dino{}
atmosphere uses a Gaussian latitude grid (the natural collocation for the
spectral transform), while the ocean and land use a regular
latitude--longitude grid. The longitudes coincide, so coupling reduces to a
\emph{one-dimensional} interpolation in latitude. \model{} uses differentiable
piecewise-linear interpolation in both directions,
\begin{equation}
  f^{\,\mathrm{Gauss}} = \mathcal{I}\big[\lat_{\mathrm{lin}}\to\lat_{\mathrm{Gauss}}\big]\,f^{\,\mathrm{lin}},
  \qquad
  f^{\,\mathrm{lin}} = \mathcal{I}\big[\lat_{\mathrm{Gauss}}\to\lat_{\mathrm{lin}}\big]\,f^{\,\mathrm{Gauss}},
  \label{eq:regrid}
\end{equation}
implemented with \code{jnp.interp} vectorized over longitude columns
(\code{jax.vmap}). Using \code{jnp.interp} rather than a looped \code{np.interp} is
what keeps the regrid differentiable, so gradients flow across the grid boundary
in both directions.
\codref{chronos\_esm/coupler/dino\_coupling.py}{\code{make\_regridders\_jax} builds $\mathcal{I}$ of Eq.~\eqref{eq:regrid}}

\begin{remark}[Linear vs.\ conservative remapping]
The latitudinal interpolation of Eq.~\eqref{eq:regrid} is \emph{not} flux-conserving:
it does not guarantee that the area integral of a flux is preserved across the
grid change. A separate dense/sparse matrix remapper exists
(\code{regrid\_field}, $\text{Dest}=W\,\text{Src}$) intended for conservative
weights, but in the current configuration it short-circuits to identity for
collocated grids. The de-meaning step of Section~\ref{sec:fluxcorr} (which
enforces zero global-mean heat and freshwater on the ocean grid) is what actually
guards global balance, not the regrid.
\codref{chronos\_esm/coupler/regrid.py}{\code{regrid\_field}, \code{Regridder} (matrix remap, identity for collocated grids)}
\end{remark}

\section{Coupling time-stepping: \texttt{step} vs.\ \texttt{step\_fast}}

A coupling interval contains many sub-steps: the atmosphere is advanced with its
ImEx integrator (Chapter~\ref{chap:ch09_atmosphere_multilevel}) over
$n_{\mathrm{atm}}$ steps, and the ocean is advanced over $n_{\mathrm{sub}}$
sub-steps via \code{jax.lax.scan}, where
\begin{equation}
  n_{\mathrm{atm}} = \lceil s_{\mathrm{atm}}\,\Delta T_{\mathrm{int}}\rceil,
  \qquad
  n_{\mathrm{sub}} = \big\lceil 86400\,\Delta T_{\mathrm{int}}/\Delta t_{\mathrm{ocean}}\big\rceil,
\end{equation}
with $\Delta T_{\mathrm{int}}$ the interval in days, $s_{\mathrm{atm}}$ the
atmosphere steps-per-day, and $\Delta t_{\mathrm{ocean}}$ the ocean time step.
\model{} exposes the same interval body through two entry points with different
performance/differentiability trade-offs:

\begin{description}
  \item[\code{step}] --- the differentiable, variable-interval path. The
        atmosphere advance and the ocean \code{scan} are wrapped in gradient
        checkpointing (\code{jax.checkpoint}, ``remat'') so the reverse-mode
        adjoint fits in memory. This is the path used for \code{jax.grad} and
        data assimilation (Chapter~\ref{chap:ch15_differentiable_applications}).
  \item[\code{step\_fast}] --- the forward-only path for long control and forced
        runs. The \emph{entire} interval (regrid, spectral transforms, atmosphere
        steps, ocean scan) is fused into one \code{jax.jit} program at a
        \emph{fixed} interval, with no remat. Removing the checkpointing lets XLA
        fuse the program, which the source reports is $\sim\!10\times$ faster than
        the eager \code{step}.
\end{description}

Both call the shared \code{\_advance} body of Eq.~\eqref{eq:cycle}; they differ
only in the \code{remat} flag and whether the interval is static.
\codref{chronos\_esm/coupler/dino\_step.py}{\code{DinoCoupledModel.step} (remat) and \code{step\_fast} (fused jit)}

\begin{remark}[Honest status]
The differentiable \code{DinoCoupledModel} is the library promotion of the
standalone control harness \code{experiments/run\_dino\_coupled.py}; it is the
\code{atmos\_backend='dino'} model. The legacy single-level coupler in
\code{main.step\_coupled} (Chapter~\ref{chap:ch10_atmosphere_singlelevel}) is
untouched and is the path the older \code{jax.grad} demos exercise. A second
honesty note inherited from the AMOC metric: on the present
diagnostic-velocity ocean, $\mathrm{d}(\text{AMOC})/\mathrm{d}(\text{density})=0$,
so a CO$_2$ or hosing perturbation entering through the heat/freshwater budget here
does not yet drive overturning gradients --- closing that is the prognostic-momentum
work of Chapter~\ref{chap:ch07_ocean_prognostic_core}.
\end{remark}

\section{Exercises}

\begin{exercise}[Restoring strength and forcing suppression]
Using Eq.~\eqref{eq:haney}, compute the Haney coefficient $\lambda$ for restoring
timescales $\tau=10,30,90,365$ days (take $\Delta z_1$ from the model's top layer,
$\densz=1025$, $c_{p,w}=3985$). For each, predict the equilibrium SST anomaly
$\Delta T = F/\lambda$ to a doubled-CO$_2$ forcing $F=3.7\ \mathrm{W\,m^{-2}}$
(Eq.~\eqref{eq:myhre}). At what $\tau$ does $\Delta T$ exceed $1\ \mathrm{K}$?
Explain why control mode is unusable for climate sensitivity and free mode
(Eq.~\eqref{eq:qflux}) is required.
\end{exercise}

\begin{exercise}[Diagnosing the implied q-flux]
Construct a \code{DinoCoupledModel} in control mode (\code{q\_flux=None},
\code{restore\_tau\_days=30}) and integrate to a quasi-steady state with
\code{step\_fast}. At each interval accumulate the restoring flux
$\lambda\,(T^{\star}-T_s)$ (Eq.~\eqref{eq:haney}) and form its time mean
$\bar{Q}$. Map $\bar{Q}$: where is the implied correction largest, and what known
SST biases (e.g.\ the tropical cold tongue) does its sign reveal? Then rebuild the
model in free mode passing this $\bar{Q}$ as \code{q\_flux} with
\code{restore\_tau\_days=3650} and verify the mean SST stays close to the control.
\end{exercise}

\begin{exercise}[Wind-stress clamp as an instability detector]
The stress components are clamped to $[-0.3,0.3]\ \mathrm{Pa}$
(Eq.~\eqref{eq:stressclamp}). Instrument \code{ocean\_fluxes\_jax} to record the
fraction of ocean points hitting the clamp each interval during a multi-year run.
Argue why a \emph{rising} clamp-hit fraction is a leading indicator of an
atmospheric instability rather than a feature of the climatology, and relate it to
the $25\ \mathrm{m\,s^{-1}}$ WISHE cap in the heat fluxes (Eq.~\eqref{eq:windeff}).
\end{exercise}

\begin{exercise}[Differentiable CO$_2$ sensitivity]
Because $F_{\mathrm{CO_2}}$ (Eq.~\eqref{eq:myhre}) and the whole coupling interval
are pure \jax{}, the map $C\mapsto \overline{\mathrm{SST}}$ after $N$ intervals is
differentiable. In free mode, use \code{jax.grad} through \code{step} to compute
$\mathrm{d}\,\overline{\mathrm{SST}}/\mathrm{d}C$ at $C=400\ \mathrm{ppm}$, and
compare it to a finite-difference estimate. Then explain, using the structure of
Section~\ref{sec:co2}, why the same gradient computed in \emph{control} mode is
much smaller, and what that implies for using gradients to calibrate climate
sensitivity.
\end{exercise}
